\chapter{End-to-end mapping: wake 경험에서 검증된 장기기억까지}
\label{bg:chapter-stc-mapping}

앞 장의 용어를 하나의 lifecycle로 조립한다. \concept{wake-state}{Wake 상태}는 사용자 요청을 처리하는
critical path에서 읽히고, query/session 범위에서 빠르게 변할 수 있다. \concept{sleep-state}{Sleep 상태}는
응답 path 밖에서 replay·학습·압축·검증되는 후보다. ``밤''이라는 시각보다 dependency와 latency boundary가
정의의 핵심이다.

\section{두 기억 destination과 hybrid promotion}

\concept{external-memory}{외부 기억}은 text, vector, graph, log처럼 model 밖에서 provenance와 함께
관리되고 query 때 검색된다. \concept{parametric-memory}{모수적 기억}은 base weights, LoRA, adapter,
fast weights처럼 forward 함수 안에서 behavior를 바꾼다.

External memory는 쓰기·삭제·scope·source attribution이 쉽고 capacity를 독립 확장할 수 있지만 retrieval
latency, context token, index drift, tool failure가 있다. Parametric memory는 retrieval 없이 behavior와
skill을 바꿀 수 있지만 update·검증·rollback이 비싸고 exact provenance와 deletion이 어렵다. 따라서
실용적 default는 external system-of-record에 먼저 저장하고, 반복 사용·안정성·scope·break-even이
충분한 pattern만 reversible parametric artifact로 승격하는 hybrid다.

\section{모든 sleep operator를 같은 schema로 기록한다}

\begin{table}[htbp]
\centering
\caption{Sleep operator의 공통 training/serving schema}
\label{tab:bg-stc-operator-schema}
\scriptsize
\begin{tabularx}{\textwidth}{p{0.11\textwidth} p{0.12\textwidth} p{0.12\textwidth} p{0.12\textwidth} X X}
\toprule
Operator & 입력 data & 변경 state & objective/signal & validation & publication/rollback \\
\midrule
summary & episode + source & text record & fidelity+compression & QA, temporal, provenance & record version/tombstone \\
re-embed & text+encoder & vector index & retrieval utility & recall, scope, latency & index epoch/rebuild \\
graph merge & facts+time & node/edge & consistency+coverage & conflict, traversal & edge version/replay log \\
replay FT & new+coreset & LoRA/adapter & task+retention loss & new/old/safety slices & detach previous adapter \\
distill & teacher outputs & student/delta & KL+task loss & fidelity, utility, tail & teacher+student pair \\
model edit & edit+neighbors & localized weights & efficacy+locality & paraphrase, specificity & inverse/version rollback \\
global refresh & curated corpus & base weights & multi-objective & full release suite & model release rollback \\
\bottomrule
\end{tabularx}
\end{table}

각 row는 반드시 teacher/source, trainable state, loss/reward, optimizer, capacity destination, validation,
version을 갖는다. Training kernel이 없는 summary/graph merge도 data transformation objective와 evaluation,
publication이 있으므로 넓은 의미의 sleep compute에 포함된다. 반대로 아무 validation 없이 log를
append하는 것은 consolidation이 아니라 storage다.

\section{Worked example: 반복되는 사용자 correction}

사용자가 여러 session에서 ``회의 요약은 표보다 세 문장 bullet을 선호한다''고 수정했다고 하자.

\begin{enumerate}
  \item \textbf{Wake capture}: correction event와 이전 response, explicit user scope, timestamp를 append한다.
  \item \textbf{Curation}: 일시적 요청인지 stable preference인지 분류하고 중복을 cluster한다.
  \item \textbf{External update}: user profile text/graph에 source-linked preference를 versioned write한다.
  \item \textbf{Benefit measurement}: retrieval이 실제 prompt token과 latency를 얼마나 늘리고, preference가
        얼마나 자주 재사용되는지 관찰한다.
  \item \textbf{Candidate promotion}: 반복 빈도와 stable period가 threshold를 넘으면 user-scoped LoRA
        training set에 positive/negative pair와 general retention replay를 넣는다.
  \item \textbf{Validation}: preference success, unrelated style, safety, other-user non-leakage, latency를 비교한다.
  \item \textbf{Publish}: external profile은 system-of-record로 유지하고 adapter pointer만 canary 승격한다.
  \item \textbf{Evict/rollback}: preference가 바뀌면 old external record를 supersede하고 adapter를 detach 또는
        새 version으로 대체한다.
\end{enumerate}

이 예에서 external과 parametric memory는 경쟁 복제품이 아니다. 하나는 evidence와 control plane,
다른 하나는 반복되는 behavior의 serving accelerator다.

\section{Break-even을 묻는 최소 식}

Parametric promotion의 one-time sleep cost를 $C_{train}+C_{validate}+C_{publish}$, query당 retrieval에서
절약되는 기대 비용을 $\Delta C_{wake}$, 예상 유효 query 수를 $N$, 추가 risk와 maintenance를
$C_{risk}+C_{maint}$라 하면 promotion 조건의 한 baseline은

\begin{equation}
N\,\Delta C_{wake} > C_{train}+C_{validate}+C_{publish}+C_{risk}+C_{maint}.
\end{equation}

정확도/utility gain도 통화 또는 constrained objective로 별도 포함한다. 이 식은 scaling law가 아니라
측정 항목을 빠뜨리지 않기 위한 accounting identity다. $N$을 과대 추정하거나 validation cost를 0으로
두면 대부분의 기억을 weights로 보내는 잘못된 결론이 나온다.

\section{운영 decision tree}

\begin{enumerate}
  \item source와 scope가 불명확하면 학습하지 말고 quarantine한다.
  \item exact delete, temporal history, citation이 중요하면 external system-of-record를 유지한다.
  \item current query만 적응하면 query/session-local state로 제한하고 reset한다.
  \item 반복 사용되는 stable pattern이며 retrieval overhead가 크면 reversible adapter 후보를 만든다.
  \item 여러 사용자에게 일반화되고 독립 evidence가 충분하면 global refresh queue로 승격한다.
  \item 어느 경우든 capacity action과 rollback artifact가 없으면 publish하지 않는다.
\end{enumerate}

\begin{keypoint}[Study Paper를 읽을 때의 질문]
새 sleep-time 연구를 볼 때 ``성능이 올랐나''에 앞서 다음을 확인한다. Wake/sleep 경계는 무엇인가?
무슨 state가 얼마나 오래 바뀌는가? Data source와 replay는 무엇인가? Old capability를 어떻게 평가했는가?
Capacity가 찼을 때 무엇을 지우는가? Candidate는 어떻게 publish·rollback되는가? 이 여섯 질문이 논문
결과를 실제 system과 memory-device workload로 번역한다.
\end{keypoint}
